% =============================
% Tractable Circuit Zoo - Succinctness Claims and Descriptions
% Auto-generated from database.json
% Generated: 2026-05-07T23:01:47.205Z
% 
% EDITING INSTRUCTIONS:
% - Claim environments are auto-generated. Do NOT edit.
% - Description environments (claimdescription) are EDITABLE.
% - Lines starting with "% [DERIVED" indicate auto-propagated edges.
% - To sync back to JSON, run: npx tsx scripts/latex-bijection.ts --to-json <this-file>
% =============================
\documentclass[11pt]{article}

% -------- Packages --------
\usepackage[margin=1in]{geometry}
\usepackage{amsmath, amssymb, amsthm}
\usepackage{mathtools}
\usepackage{enumitem}
\usepackage{hyperref}
\usepackage{cleveref}
\usepackage{xcolor}
\usepackage{natbib}

% -------- Hyperref setup --------
\hypersetup{
  colorlinks=true,
  linkcolor=blue,
  citecolor=blue,
  urlcolor=blue
}

% -------- Theorem styles --------
\theoremstyle{plain}
\newtheorem{theorem}{Theorem}[section]
\newtheorem{lemma}[theorem]{Lemma}
\newtheorem{proposition}[theorem]{Proposition}
\newtheorem{corollary}[theorem]{Corollary}
\newtheorem{claim}{Claim}[section]

\theoremstyle{definition}
\newtheorem{definition}[theorem]{Definition}
\newtheorem{example}[theorem]{Example}

\theoremstyle{remark}
\newtheorem{remark}[theorem]{Remark}

% -------- Description environment (just indented text, no prefix) --------
\newenvironment{claimdescription}{%
  \par\noindent\ignorespaces
}{\par}

% -------- Handy macros --------
\newcommand{\R}{\mathbb{R}}
\newcommand{\N}{\mathbb{N}}
\newcommand{\eps}{\varepsilon}
% Cross-reference commands (rendered as links in the web UI)
\newcommand{\langref}[1]{\textbf{#1}}
\newcommand{\langfam}[2]{\textbf{#1$_{#2}$}}
\newcommand{\edgeref}[2]{#1 compiles to #2}
\newcommand{\nedgeref}[2]{#1 cannot compile to #2}
\newcommand{\opref}[2]{#1 supports #2}
\newcommand{\nopref}[2]{#2 is unsupported by #1}

% -------- Title info --------
\title{Tractable Circuit Zoo: Succinctness Claims}
\date{\today}

\begin{document}
\maketitle

\tableofcontents
\newpage

% =============================
\section{A. Amarilli and F. Capelli and M. Monet and P. Senellart, "Connecting Knowledge ...}
% Reference ID: Amarilli_2018
% =============================
\begin{claim}
\langref{d-SDNNF} is polynomial-time compilable to \langref{d-DNNF} \citet{Amarilli_2018}
\end{claim}
\begin{claimdescription}
As defined by \citet{Amarilli_2018}, a \langref{d-SDNNF} is simply a structured \langref{d-DNNF}, so discarding the vtree completes the compilation.
\end{claimdescription}

\begin{claim}
\langref{d-SDNNF} is polynomial-time compilable to \langref{SDNNF} \citet{Amarilli_2018}
\end{claim}
\begin{claimdescription}
As defined by \citet{Amarilli_2018}, \langref{d-SDNNF} is exactly the subset of \langref{SDNNF} satisfying determinism, so compilation is trivial.
\end{claimdescription}

\begin{claim}
\langref{d-SDNNF} has unknown polynomial-time (but has quasi-polynomial-time) compilation to \langref{uOBDD} \citet{Amarilli_2018,Bollig_2018,Beame_2015}
\end{claim}
\begin{claimdescription}
Theorem 2 of \citet{Bollig_2018} shows that \langref{SDNNF} can be compiled into \langref{nOBDD} by adapting the compilation of \citet{Beame_2015} from \langref{DNNF} to \langref{nFBDD}; proposition 2 of \citet{Bollig_2018} shows that this \langref{nOBDD} is unambiguous if the \langref{SDNNF} is deterministic.
\end{claimdescription}

\begin{claim}
\langfam{d-SDNNF}{T} is not polynomial-time compilable to \langref{nFBDD} \citet{Amarilli_2018}
\end{claim}
\begin{claimdescription}
This follows because \langfam{dec-SDNNF}{T} is a polynomial-time sublanguage of \langfam{d-SDNNF}{T} and \langfam{dec-SDNNF}{T} does not compile to \langref{nFBDD} in polynomial time.
\end{claimdescription}

\begin{claim}
\langfam{d-SDNNF}{T} is polynomial-time compilable to \langfam{SDNNF}{T} \citet{Amarilli_2018}
\end{claim}
\begin{claimdescription}
For any fixed vtree $T$, \langfam{d-SDNNF}{T} is exactly the subset of \langfam{SDNNF}{T} satisfying determinism, so compilation is trivial.
\end{claimdescription}

\begin{claim}
\langfam{d-SDNNF}{T} has unknown polynomial-time (but has quasi-polynomial-time) compilation to \langfam{uOBDD}{<} \citet{Amarilli_2018,Bollig_2018,Beame_2015}
\end{claim}
\begin{claimdescription}
Apply the common-order simulation of quasi-note.tex to \langfam{d-SDNNF}{T}. The output order depends only on $T$, and the Bollig--Buttkus certificate argument shows that deterministic inputs yield unambiguous nondeterministic OBDDs. The polynomial lower bound follows from \langfam{d-SDNNF}{T} not compiling to \langref{nFBDD} in polynomial time and the polynomial inclusion \langfam{uOBDD}{<} $\subseteq$ \langref{nFBDD}.
\end{claimdescription}

\begin{claim}
\langref{dec-DNNF} is polynomial-time compilable to \langref{d-DNNF} \citet{Amarilli_2018}
\end{claim}
\begin{claimdescription}
As defined by \citet{Amarilli_2018}, \langref{dec-DNNF} is the subset of \langref{d-DNNF} satisfying the decision property.
\end{claimdescription}

\begin{claim}
\langref{dec-DNNF} has unknown polynomial-time (but has quasi-polynomial-time) compilation to \langref{FBDD} \citet{Beame_2013}
\end{claim}
\begin{claimdescription}
Corollary 3.2 of \citet{Beame_2013} show quasi-polynomial compilation from \langref{dec-DNNF} to \langref{FBDD}
\end{claimdescription}

\begin{claim}
\langref{dec-SDNNF} is polynomial-time compilable to \langref{d-SDNNF} \citet{Amarilli_2018}
\end{claim}
\begin{claimdescription}
As defined by \citet{Amarilli_2018}, \langref{dec-SDNNF} is the subset of \langref{d-SDNNF} satisfying the decision property, so compilation is trivial.
\end{claimdescription}

\begin{claim}
\langref{dec-SDNNF} is polynomial-time compilable to \langref{dec-DNNF} \citet{Amarilli_2018}
\end{claim}
\begin{claimdescription}
As defined by \citet{Amarilli_2018}, a \langref{dec-SDNNF} is simply a structured \langref{dec-DNNF}, so discarding the v-tree completes the compilation.
\end{claimdescription}

\begin{claim}
\langref{dec-SDNNF} has unknown polynomial-time (but has quasi-polynomial-time) compilation to \langref{OBDD} \citet{Amarilli_2018,Bollig_2018,Beame_2015}
\end{claim}
\begin{claimdescription}
Theorem 2 of \citet{Bollig_2018} shows that \langref{SDNNF} can be compiled into \langref{nOBDD} by adapting the compilation of \citet{Beame_2015} from \langref{DNNF} to \langref{nFBDD}; proposition 2 of \citet{Bollig_2018} shows that this \langref{nOBDD} is unambiguous if the \langref{SDNNF} is deterministic. Additionally, the same compilation yields an \langref{OBDD} if the input is a \langref{dec-SDNNF} since there are no $\vee$-gates that are not decision gates in a \langref{dec-SDNNF}.
\end{claimdescription}

\begin{claim}
\langfam{dec-SDNNF}{T} is polynomial-time compilable to \langfam{d-SDNNF}{T} \citet{Amarilli_2018}
\end{claim}
\begin{claimdescription}
For any fixed vtree $T$, \langfam{dec-SDNNF}{T} is the subset of \langfam{d-SDNNF}{T} satisfying the decision property, so compilation is trivial.
\end{claimdescription}

\begin{claim}
\langfam{dec-SDNNF}{T} is not polynomial-time compilable to \langref{nFBDD} \citet{Amarilli_2018}
\end{claim}
\begin{claimdescription}
This is the fixed-vtree form of Proposition 3.10 of \citet{Amarilli_2018}.
\end{claimdescription}

\begin{claim}
\langfam{dec-SDNNF}{T} has unknown polynomial-time (but has quasi-polynomial-time) compilation to \langfam{OBDD}{<} \citet{Amarilli_2018,Bollig_2018,Beame_2015}
\end{claim}
\begin{claimdescription}
The common-order simulation of quasi-note.tex specializes to \langfam{dec-SDNNF}{T}: for every fixed source vtree $T$, it chooses one output order $<$ depending only on $T$ and produces a quasipolynomial-size \langfam{OBDD}{<}. The polynomial lower bound follows from \langfam{dec-SDNNF}{T} not compiling to \langref{FBDD} in polynomial time and the polynomial inclusion \langfam{OBDD}{<} $\subseteq$ \langref{FBDD}.
\end{claimdescription}

\begin{claim}
\langref{DNF} is not quasi-polynomial-time compilable to \langref{d-SDNNF} \citet{Amarilli_2018}
\end{claim}
\begin{claimdescription}
Corollary 7.6 of \citet{Amarilli_2018} gives a lower bound of $2^{\Omega(tw(\phi))}$ on the smallest \langref{d-SDNNF} computing any monotone \langref{DNF} of constant arity and degree. In particular, DNF families whose treewidth grows linearly with their size force exponential-size \langref{d-SDNNF} representations.
\end{claimdescription}

\begin{claim}
\langref{DNF} is polynomial-time compilable to \langfam{nOBDD}{<} \citet{Amarilli_2018,Wegener_2000}
\end{claim}
\begin{claimdescription}
Fix any variable order $<$. Each term of a \langref{DNF} can be represented as an \langfam{OBDD}{<} of linear size (a chain of decision nodes for each literal in the term). Their nondeterministic disjunction yields an \langfam{nOBDD}{<} of size $O(|\phi|)$ because every branch follows the same fixed order.
\end{claimdescription}

\begin{claim}
\langref{DNNF} has unknown polynomial-time (but has quasi-polynomial-time) compilation to \langref{nFBDD} \citet{Bodlaender_1993}
\end{claim}
\begin{claimdescription}
Section 5 of \citet{Bodlaender_1993} extends the compilation of \langref{dec-DNNF}s into \langref{FBDD}s (shown in \citet{Beame_2013}) to show that \langref{DNNF} can be quasi-polynomial-time compiled into \langref{nFBDD}.
\end{claimdescription}

\begin{claim}
\langref{FBDD} is polynomial-time compilable to \langref{dec-DNNF} \citet{Amarilli_2018}
\end{claim}
\begin{claimdescription}
Rewrite each decision node labeled $x$ as $(x \wedge D_0) \vee (\neg x \wedge D_1)$, where $D_0$ and $D_1$ are the rewritings of the 0- and 1-successors. The introduced $\vee$-gates are decision gates and the $\wedge$-gates are decomposable, so all $\vee$-gates arise from decision node rewrites, yielding a \langref{dec-DNNF} in linear time \citet{Amarilli_2018}.
\end{claimdescription}

\begin{claim}
\langref{FBDD} is polynomial-time compilable to \langref{uFBDD} \citet{Amarilli_2018}
\end{claim}
\begin{claimdescription}
As defined by \citet{Amarilli_2018}, \langref{FBDD} is the subset of \langref{uFBDD} with no $\vee$-nodes.
\end{claimdescription}

% separators=Sauerhoff
\begin{claim}
\langref{nFBDD} is not quasi-polynomial-time compilable to \langref{d-DNNF} \citet{Amarilli_2018}
\end{claim}
\begin{claimdescription}
The Sauerhoff function has an \langref{nFBDD} of size $O(n^2)$ over $n^2$ variables, but as a \langref{d-DNNF} has size $2^{\Omega(n)}$.
\end{claimdescription}

\begin{claim}
\langref{nFBDD} is polynomial-time compilable to \langref{DNNF} \citet{Amarilli_2018}
\end{claim}
\begin{claimdescription}
Rewrite each decision node labeled $x$ as $(x \wedge D_0) \vee (\neg x \wedge D_1)$, where $D_0, D_1$ are the rewritings of the 0- and 1-successors. The $\wedge$-gates are decomposable, yielding a \langref{DNNF} in linear time \citet{Amarilli_2018}.
\end{claimdescription}

\begin{claim}
\langref{nFBDD} is not quasi-polynomial-time compilable to \langref{FBDD} \citet{Amarilli_2018}
\end{claim}
\begin{claimdescription}
Consider the Boolean functions $\psi_n,\phi_n:\{0,1\}^{n^2}\to\{0,1\}$ that respectively test whether, in an $n \times n$ Boolean matrix, either the number of 1's is odd and there is a full row of 1's, or the number of 1's is even and there is a column full of 1's. An \langref{FBDD} of size $O(n^2)$ may simulate either $\psi_n$ or $\phi_n$, and thus the \langref{uFBDD} $\psi_n\vee\phi_n$ is of size $O(n^2)$; however, an \langref{FBDD} for $\psi_n\vee\phi_n$ is provably of size $2^{\Omega(n^{1/2})}$ \citet{Amarilli_2018}.
\end{claimdescription}

% separators=Sauerhoff
\begin{claim}
\langref{nFBDD} is not quasi-polynomial-time compilable to \langref{uFBDD} \citet{Amarilli_2018}
\end{claim}
\begin{claimdescription}
The Sauerhoff function has an \langref{nFBDD} of size $O(n^2)$ over $n^2$ variables, but as a \langref{d-DNNF}, which generalize \langref{uFBDD}, has size $2^{\Omega(n)}$.
\end{claimdescription}

\begin{claim}
\langref{nOBDD} is polynomial-time compilable to \langref{nFBDD} \citet{Amarilli_2018}
\end{claim}
\begin{claimdescription}
As defined by \citet{Amarilli_2018}, \langref{nOBDD} is the subset of \langref{nFBDD} satisfying an order of the variables, so compilation entails discarding the order.
\end{claimdescription}

\begin{claim}
\langref{nOBDD} is polynomial-time compilable to \langref{SDNNF} \citet{Amarilli_2018}
\end{claim}
\begin{claimdescription}
Construct a right-linear v-tree from the variable order, then rewrite each decision node labeled $x$ as $(x \wedge D_0) \vee (\neg x \wedge D_1)$, where $D_0, D_1$ are the rewritings of the 0- and 1-successors. The v-tree alignment ensures every $\wedge$-gate respects the variable partition, yielding an \langref{SDNNF} in linear time \citet{Amarilli_2018}.
\end{claimdescription}

\begin{claim}
\langfam{nOBDD}{<} is polynomial-time compilable to \langfam{SDNNF}{T} \citet{Amarilli_2018}
\end{claim}
\begin{claimdescription}
Fix any variable order $<$ and let $T$ be the corresponding right-linear vtree. Rewriting each decision node labeled $x$ as $(x \wedge D_0) \vee (\neg x \wedge D_1)$ yields an output respecting $T$, so \langfam{nOBDD}{<} compiles to \langfam{SDNNF}{T} in linear time.
\end{claimdescription}

\begin{claim}
\langref{OBDD} is polynomial-time compilable to \langref{dec-SDNNF} \citet{Amarilli_2018}
\end{claim}
\begin{claimdescription}
Construct a right-linear v-tree from the variable order, then rewrite each decision node labeled $x$ as $(x \wedge D_0) \vee (\neg x \wedge D_1)$, where $D_0, D_1$ are the rewritings of the 0- and 1-successors. The v-tree alignment ensures structuredness, and since all $\vee$-gates arise from decision node rewrites, every $\vee$-gate is a decision gate, yielding a \langref{dec-SDNNF} in linear time \citet{Amarilli_2018}.
\end{claimdescription}

\begin{claim}
\langref{OBDD} is polynomial-time compilable to \langref{uOBDD} \citet{Amarilli_2018}
\end{claim}
\begin{claimdescription}
As defined by \citet{Amarilli_2018}, \langref{OBDD} is the subset of \langref{uOBDD} without any $\vee$-node.
\end{claimdescription}

\begin{claim}
\langfam{OBDD}{<} is polynomial-time compilable to \langfam{dec-SDNNF}{T} \citet{Amarilli_2018}
\end{claim}
\begin{claimdescription}
Fix any variable order $<$ and let $T$ be the corresponding right-linear vtree. Rewriting each decision node labeled $x$ as $(x \wedge D_0) \vee (\neg x \wedge D_1)$ yields a structured output that respects $T$, and all introduced $\vee$-gates are decision gates, producing \langfam{dec-SDNNF}{T} in linear time.
\end{claimdescription}

\begin{claim}
\langfam{OBDD}{<} is polynomial-time compilable to \langfam{uOBDD}{<} \citet{Amarilli_2018,Wegener_2000}
\end{claim}
\begin{claimdescription}
For any fixed variable order $<$, every \langfam{OBDD}{<} is already an unambiguous nondeterministic OBDD respecting the same order. Thus \langfam{OBDD}{<} is a subset of \langfam{uOBDD}{<}, and the compilation is the identity.
\end{claimdescription}

\begin{claim}
\langref{SDNNF} is polynomial-time compilable to \langref{DNNF} \citet{Amarilli_2018}
\end{claim}
\begin{claimdescription}
As defined by \citet{Amarilli_2018}, a structured \langref{DNNF} is a triple of the corresponding \langref{DNNF}, the v-tree $T$, and a mapping $\rho$; as such, extracting a \langref{DNNF} from a \langref{SDNNF} is trivial.
\end{claimdescription}

\begin{claim}
\langref{SDNNF} has unknown polynomial-time (but has quasi-polynomial-time) compilation to \langref{nOBDD} \citet{Bollig_2018,Beame_2015}
\end{claim}
\begin{claimdescription}
Theorem 2 of \citet{Bollig_2018} shows that \langref{SDNNF} can be compiled into \langref{nOBDD} by adapting the compilation of \citet{Beame_2015} from \langref{DNNF} to \langref{nFBDD}.
\end{claimdescription}

\begin{claim}
\langfam{SDNNF}{T} has unknown polynomial-time (but has quasi-polynomial-time) compilation to \langfam{nOBDD}{<} \citet{Amarilli_2018,Bollig_2018,Beame_2015}
\end{claim}
\begin{claimdescription}
The Bollig--Buttkus simulation of structured DNNFs by nondeterministic OBDDs can use the vtree leaf-count orientation of quasi-note.tex. For every fixed source vtree $T$, this chooses one output order $<$ depending only on $T$, not on the input circuit, and gives a quasipolynomial-size \langfam{nOBDD}{<}. The polynomial lower bound follows from \langfam{d-SDNNF}{T} not compiling to \langref{nFBDD} in polynomial time, together with the polynomial inclusions \langfam{d-SDNNF}{T} $\subseteq$ \langfam{SDNNF}{T} and \langfam{nOBDD}{<} $\subseteq$ \langref{nFBDD}.
\end{claimdescription}

\begin{claim}
\langref{uFBDD} is polynomial-time compilable to \langref{d-DNNF} \citet{Amarilli_2018}
\end{claim}
\begin{claimdescription}
Rewrite each decision node labeled $x$ as $(x \wedge D_0) \vee (\neg x \wedge D_1)$, where $D_0, D_1$ are the rewritings of the 0- and 1-successors. The $\wedge$-gates are decomposable, and unambiguity of the \langref{uFBDD} ensures all introduced $\vee$-gates are deterministic, yielding a \langref{d-DNNF} in linear time \citet{Amarilli_2018}.
\end{claimdescription}

\begin{claim}
\langref{uFBDD} is not quasi-polynomial-time compilable to \langref{FBDD} \citet{Bova_2016,Amarilli_2018,Wegener_2000}
\end{claim}
\begin{claimdescription}
Consider the Boolean functions $\phi_n, \psi_n: \{0,1\}^{n^2}\to\{0,1\}$, which respectively test whether, in an $n\times n$ Boolean matrix, that the number of 1's is odd and there is a row full of 1's, and that the number of 1's is even and there is a column full of 1's. \citet{Bova_2016} and \citet{Wegener_2000} show that an \langref{FBDD} expression for $\phi_n \vee \psi_n$ has necessarily size $2^{\Omega(n^{1/2})}$; but since an \langref{FBDD} of size $O(n^2)$ may test $\phi_n$ or $\psi_n$ separately (shown by Proposition 3.1 of \citet{Amarilli_2018}), a \langref{uFBDD} for $\phi_n\vee \psi_n$ may be obtained by simply adding an $\vee$-gate joining these two \langref{FBDD} expressions, since $\phi_n \wedge \psi_n$ is unsatisfiable.
\end{claimdescription}

\begin{claim}
\langref{uFBDD} is polynomial-time compilable to \langref{nFBDD} \citet{Amarilli_2018}
\end{claim}
\begin{claimdescription}
As defined by \citet{Amarilli_2018}, \langref{uFBDD} is the subset of \langref{nFBDD} satisfying unambiguity.
\end{claimdescription}

\begin{claim}
\langref{uOBDD} is polynomial-time compilable to \langref{d-SDNNF} \citet{Amarilli_2018}
\end{claim}
\begin{claimdescription}
Construct a right-linear v-tree from the variable order, then rewrite each decision node labeled $x$ as $(x \wedge D_0) \vee (\neg x \wedge D_1)$, where $D_0, D_1$ are the rewritings of the 0- and 1-successors. The v-tree alignment ensures structuredness, and unambiguity of the \langref{uOBDD} makes all introduced $\vee$-gates deterministic, yielding a \langref{d-SDNNF} in linear time \citet{Amarilli_2018}.
\end{claimdescription}

\begin{claim}
\langref{uOBDD} is polynomial-time compilable to \langref{nOBDD} \citet{Amarilli_2018}
\end{claim}
\begin{claimdescription}
As defined by \citet{Amarilli_2018}, \langref{uOBDD} is the subset of \langref{nOBDD} satisfying unambiguity.
\end{claimdescription}

\begin{claim}
\langref{uOBDD} is polynomial-time compilable to \langref{uFBDD} \citet{Amarilli_2018}
\end{claim}
\begin{claimdescription}
As defined by \citet{Amarilli_2018}, \langref{uOBDD} is the subset of \langref{uFBDD} satisfying an order of the variables, so compilation entails discarding the order.
\end{claimdescription}

\begin{claim}
\langfam{uOBDD}{<} is polynomial-time compilable to \langfam{d-SDNNF}{T} \citet{Amarilli_2018}
\end{claim}
\begin{claimdescription}
Fix any variable order $<$ and let $T$ be the corresponding right-linear vtree. Rewriting each decision node labeled $x$ as $(x \wedge D_0) \vee (\neg x \wedge D_1)$ yields a structured output respecting $T$, and unambiguity ensures determinism of introduced $\vee$-gates, giving \langfam{d-SDNNF}{T} in linear time.
\end{claimdescription}

\begin{claim}
\langfam{uOBDD}{<} is polynomial-time compilable to \langfam{nOBDD}{<} \citet{Amarilli_2018,Wegener_2000}
\end{claim}
\begin{claimdescription}
For any fixed variable order $<$, every unambiguous nondeterministic OBDD respecting $<$ is a nondeterministic OBDD respecting the same order. Thus \langfam{uOBDD}{<} is a subset of \langfam{nOBDD}{<}, and the compilation is the identity.
\end{claimdescription}


% =============================
\section{A. Darwiche and P. Marquis, "A Knowledge Compilation Map", Journal of Artificial...}
% Reference ID: Darwiche_2002
% =============================
\begin{claim}
\langref{CNF} is polynomial-time compilable to \langref{NNF} \citet{Darwiche_2002}
\end{claim}
\begin{claimdescription}
As defined by \citet{Darwiche_2002}, \langref{CNF} is a subset of \langref{NNF}. Therefore, a polynomial compilation is trivial.
\end{claimdescription}

% separators=Parity
\begin{claim}
\langref{CNF} is not quasi-polynomial-time compilable to \langref{PI} \citet{Darwiche_2002}
\end{claim}
\begin{claimdescription}
It is well-known that some \langref{CNF} formulas (i.e. the negation of the parity function $\neg PARITY_n=\neg\bigoplus_{i=1}^n x_i$) have exponentially many prime implicants.
\end{claimdescription}

\begin{claim}
\langref{d-DNNF} is polynomial-time compilable to \langref{DNNF} \citet{Darwiche_2002}
\end{claim}
\begin{claimdescription}
As formalized by \citet{Darwiche_2002} and originally introduced in \citet{Darwiche_2001a}, \langref{d-DNNF} is a subset of \langref{DNNF} satisfying the determinism property, making a polynomial compilation trivial.
\end{claimdescription}

% separators=$OR_n$
\begin{claim}
\langref{d-DNNF} is not quasi-polynomial-time compilable to \langref{MODS} \citet{Darwiche_2002}
\end{claim}
\begin{claimdescription}
The $OR_n$ function $\bigvee_{i=1}^nx_i$ can be represented in $O(n)$ as a \langref{d-DNNF} formula but has $2^n-1$ models.
\end{claimdescription}

\begin{claim}
\langref{DNF} is not polynomial-time compilable to \langref{d-DNNF} assuming the polynomial hierarchy does not collapse \citet{Darwiche_2002}
\end{claim}
\begin{claimdescription}
Because counting the models of a \langref{d-DNNF} is in P and counting the models of a \langref{DNF} is \#P-complete, \langref{DNF} cannot be polynomially compiled into \langref{d-DNNF} unconditionally. Such a compilation implies P = \#P, collapsing the polynomial hierarchy \citep{Darwiche_2002}.
\end{claimdescription}

\begin{claim}
\langref{DNF} is polynomial-time compilable to \langref{DNNF} \citet{Darwiche_2002}
\end{claim}
\begin{claimdescription}
As defined by \citet{Darwiche_2002}, \langref{DNF} is a subset of \langref{DNNF}. Therefore, a polynomial compilation is trivial.
\end{claimdescription}

% separators=Parity
\begin{claim}
\langref{DNF} is not quasi-polynomial-time compilable to \langref{IP} \citet{Darwiche_2002}
\end{claim}
\begin{claimdescription}
It is well-known that some \langref{DNF} formulas (i.e. the parity function $PARITY_n=\bigoplus_{i=1}^n x_i$) have exponentially many prime implicates.
\end{claimdescription}

\begin{claim}
\langref{DNNF} is polynomial-time compilable to \langref{NNF} \citet{Darwiche_2002}
\end{claim}
\begin{claimdescription}
As defined by \citet{Darwiche_2002}, \langref{DNNF} is a subset of \langref{NNF}, so a polynomial compilation is trivial.
\end{claimdescription}

\begin{claim}
\langref{FBDD} is polynomial-time compilable to \langref{d-DNNF} \citet{Darwiche_2002}
\end{claim}
\begin{claimdescription}
As defined by \citet{Darwiche_2002}, \langref{FBDD} is a subset of \langref{d-DNNF}, so a polynomial compilation is trivial.
\end{claimdescription}

% separators=Pairwise DNF
\begin{claim}
\langref{IP} is not quasi-polynomial-time compilable to \langref{CNF} \citet{Darwiche_2002}
\end{claim}
\begin{claimdescription}
Consider the \langref{DNF} formula $\Pi_n=\bigvee_{i=0}^{n-1}(x_{2i} \wedge x_{2i+1})$. This formula is in prime implicants form and each clause in $\Pi_n$ is an essential prime implicant of it. Hence its negation $\neg\Pi_n\in IP$.

Since the early work of Quine, we know that the number of essential prime implicates of a formula is a lower bound on the number of clauses that can be found in any \langref{CNF} representation of it. However, since $\Pi_n$ has $2^n$ prime implicants, $\neg\Pi_n$ has $2^n$ prime implicates, so any \langref{CNF} representation of $\neg\Pi_n$ must have at least $2^n$ clauses.
\end{claimdescription}

\begin{claim}
\langref{IP} is polynomial-time compilable to \langref{DNF} \citet{Darwiche_2002}
\end{claim}
\begin{claimdescription}
As defined by \citet{Darwiche_2002}, \langref{IP} is a subset of \langref{DNF}. Therefore, a polynomial compilation is trivial.
\end{claimdescription}

% separators=Clique
\begin{claim}
\langref{IP} is not quasi-polynomial-time compilable to \langref{FBDD} \citet{Darwiche_2002,Wegener_1987}
\end{claim}
\begin{claimdescription}
The negation of the clique detection function $\neg\Sigma_{n,k}$ has $O(n^3)$ prime implicants (hence polynomial \langref{IP} size), but any \langref{FBDD} for $\neg\Sigma_{n,k}$ has exponential size for suitable $k$ \citet{Darwiche_2002,Wegener_1987}.
\end{claimdescription}

% separators=$OR_n$
\begin{claim}
\langref{IP} is not quasi-polynomial-time compilable to \langref{MODS} \citet{Darwiche_2002}
\end{claim}
\begin{claimdescription}
The $OR_n$ function $\bigvee_{i=1}^nx_i$ can be represented in $O(n)$ as a \langref{IP} formula but has $2^n-1$ models.
\end{claimdescription}

\begin{claim}
\langref{MODS} is polynomial-time compilable to \langref{CNF} \citet{Darwiche_2002}
\end{claim}
\begin{claimdescription}
A Shannon tree for any \langref{MODS} formula can be generated in polynomial time. Paths from the root to 0-leaves yield a \langref{CNF} representation \citet{Darwiche_2002}.
\end{claimdescription}

\begin{claim}
\langref{MODS} is polynomial-time compilable to \langref{d-DNNF} \citet{Darwiche_2002}
\end{claim}
\begin{claimdescription}
As defined by \citet{Darwiche_2002}, \langref{MODS} is a subset of \langref{d-DNNF}. Therefore, a polynomial compilation is trivial.
\end{claimdescription}

\begin{claim}
\langref{MODS} is polynomial-time compilable to \langref{DNF} \citet{Darwiche_2002}
\end{claim}
\begin{claimdescription}
As defined by \citet{Darwiche_2002}, \langref{MODS} is a subset of \langref{DNF} (any \langref{MODS} formula can be written as a disjunction of the individual models). Therefore, a polynomial compilation is trivial.
\end{claimdescription}

\begin{claim}
\langref{MODS} is polynomial-time compilable to \langfam{OBDD}{<} \citet{Darwiche_2002}
\end{claim}
\begin{claimdescription}
A Shannon tree for any \langref{MODS} formula following order $<$ can be generated in polynomial time. This tree can be reduced into a corresponding \langfam{OBDD}{<} in polynomial time \citet{Darwiche_2002}.
\end{claimdescription}

\begin{claim}
\langref{OBDD} is polynomial-time compilable to \langref{FBDD} \citet{Darwiche_2002}
\end{claim}
\begin{claimdescription}
As defined by \citet{Darwiche_2002}, \langref{OBDD} is a subset of \langref{FBDD}. Therefore a polynomial time compilation is trivial.
\end{claimdescription}

% separators=Pairwise equivalence
\begin{claim}
\langref{OBDD} is not quasi-polynomial-time compilable to \langfam{OBDD}{<} \citet{Darwiche_2002}
\end{claim}
\begin{claimdescription}
The function $\bigwedge_{i=1}^n(x_i\Leftrightarrow y_i)$ has an \langfam{OBDD}{<} representation of size polynomial in $n$ whenever $<$ satisfies $x_1<y_1<...<x_n<y_n$ but exponential in $n$ if $<$ satisfies $x_1<x_2<...<x_n<y_1<y_2<...<y_n$
\end{claimdescription}

% separators=Parity
\begin{claim}
\langfam{OBDD}{<} is not quasi-polynomial-time compilable to \langref{CNF} \citet{Darwiche_2002}
\end{claim}
\begin{claimdescription}
The parity function $PARITY_n=\bigoplus_{i=1}^n x_i$ has linear size \langfam{OBDD}{<} representations but only exponential size \langref{CNF} representations.
\end{claimdescription}

% separators=Parity
\begin{claim}
\langfam{OBDD}{<} is not quasi-polynomial-time compilable to \langref{DNF} \citet{Darwiche_2002}
\end{claim}
\begin{claimdescription}
The parity function $PARITY_n=\bigoplus_{i=1}^n x_i$ has linear size \langfam{OBDD}{<} representations but only exponential size \langref{DNF} representations.
\end{claimdescription}

\begin{claim}
\langfam{OBDD}{<} is polynomial-time compilable to \langref{OBDD} \citet{Darwiche_2002}
\end{claim}
\begin{claimdescription}
As defined by \citet{Darwiche_2002}, \langfam{OBDD}{<} is a subset of \langref{OBDD}. Therefore, a polynomial compilation is trivial.
\end{claimdescription}

\begin{claim}
\langref{PI} is polynomial-time compilable to \langref{CNF} \citet{Darwiche_2002}
\end{claim}
\begin{claimdescription}
As defined by \citet{Darwiche_2002}, \langref{PI} is a subset of \langref{CNF}. Therefore, a polynomial compilation is trivial.
\end{claimdescription}

% separators=Pairwise CNF
\begin{claim}
\langref{PI} is not quasi-polynomial-time compilable to \langref{DNF} \citet{Darwiche_2002}
\end{claim}
\begin{claimdescription}
The pairwise CNF $\Sigma_n = \bigwedge_{i=0}^{n-1}(x_{2i} \vee x_{2i+1})$ has $2^n$ prime implicates; its negation $\neg\Sigma_n \in PI$ therefore has $2^n$ essential prime implicants, forcing any \langref{DNF} representation to have at least $2^n$ terms \citet{Darwiche_2002}.
\end{claimdescription}

% separators=Clique
\begin{claim}
\langref{PI} is not quasi-polynomial-time compilable to \langref{FBDD} \citet{Wegener_1987,Darwiche_2002}
\end{claim}
\begin{claimdescription}
The clique detection function $\Sigma_{n,k}$ has $O(n^3)$ prime implicates (hence polynomial \langref{PI} size), but any \langref{FBDD} for $\Sigma_{n,k}$ has exponential size for suitable $k$ \citet{Wegener_1987,Darwiche_2002}.
\end{claimdescription}

% separators=$OR_n$
\begin{claim}
\langref{PI} is not quasi-polynomial-time compilable to \langref{MODS} \citet{Darwiche_2002}
\end{claim}
\begin{claimdescription}
The $OR_n$ function $\bigvee_{i=1}^nx_i$ can be represented in $O(n)$ as a \langref{PI} formula but has $2^n-1$ models.
\end{claimdescription}


% =============================
\section{I. Wegener, "Branching Programs and Binary Decision Diagrams: Theory and Applica...}
% Reference ID: Wegener_2000
% =============================
% separators=Generalized_HWB
\begin{claim}
\langfam{cSDD}{T} is not quasi-polynomial-time compilable to \langfam{OBDD}{<} \citet{Bova_2016_a,Bryant_1986,Wegener_2000}
\end{claim}
\begin{claimdescription}
The generalized hidden weighted bit functions $F_n$ of \citet{Bova_2016_a} have compressed SDDs of size $O(n^3)$ respecting the explicit vtrees used in the construction. Those finite vtrees can be packed, after renaming variables into disjoint blocks, into one global source vtree. For the lower bound, conditioning the auxiliary variables of $F_n$ recovers $HWB_n$; conditioning an \langfam{OBDD}{<} preserves the same variable order and cannot increase size. Since Bryant's lower bound gives exponential OBDD size for $HWB_n$ for every variable order \citet{Bryant_1986,Wegener_2000}, every fixed-order OBDD for $F_n$ is exponential. Hence \langfam{cSDD}{T} is not quasi-polynomial-time compilable to \langfam{OBDD}{<}.
\end{claimdescription}

\begin{claim}
\langfam{nOBDD}{<} is polynomial-time compilable to \langref{nOBDD} \citet{Wegener_2000}
\end{claim}
\begin{claimdescription}
\langref{nOBDD} is the union of fixed-order languages \langfam{nOBDD}{<}. Hence any \langfam{nOBDD}{<} formula is already an \langref{nOBDD} formula, so the compilation is trivial and polynomial-time.
\end{claimdescription}

% separators=HWB
\begin{claim}
\langref{uOBDD} is not quasi-polynomial-time compilable to \langref{OBDD} \citet{Bryant_1986,Wegener_2000}
\end{claim}
\begin{claimdescription}
The hidden weighted bit function $HWB_n$ has polynomial-size unambiguous nondeterministic OBDDs under a fixed order: $HWB_n = \bigvee_{k=1}^n(E_k^n(x) \wedge x_k)$, where each branch for $k$ tests the variables in the same order to verify the exact Hamming weight and $x_k=1$, and the branches are disjoint. Thus the representation is unambiguous and has polynomial size, as in the 2-OBDD/FBDD constructions discussed by \citet{Wegener_2000}. Bryant's hidden-weighted-bit lower bound says every OBDD for $HWB_n$, over every variable order, has exponential size \citet{Bryant_1986,Wegener_2000}. Hence no quasi-polynomial compilation to OBDD is possible.
\end{claimdescription}

% separators=HWB
\begin{claim}
\langfam{uOBDD}{<} is not quasi-polynomial-time compilable to \langfam{OBDD}{<} \citet{Bryant_1986,Wegener_2000}
\end{claim}
\begin{claimdescription}
Use the same $HWB_n$ family as for \langref{uOBDD} versus \langref{OBDD}. The polynomial unambiguous construction can be taken to respect one fixed order for each $n$, and the finite hard instances can be renamed into disjoint variable blocks of one global order. On the target side, Bryant's lower bound is order-independent: every OBDD for $HWB_n$, and hence every \langfam{OBDD}{<} for every fixed order $<$, has exponential size \citet{Bryant_1986,Wegener_2000}. Therefore \langfam{uOBDD}{<} is not quasi-polynomial-time compilable to \langfam{OBDD}{<}.
\end{claimdescription}

\begin{claim}
\langfam{uOBDD}{<} is polynomial-time compilable to \langref{uOBDD} \citet{Wegener_2000}
\end{claim}
\begin{claimdescription}
\langref{uOBDD} is the union of fixed-order languages \langfam{uOBDD}{<}. Hence any \langfam{uOBDD}{<} formula is already a \langref{uOBDD} formula, so the compilation is trivial and polynomial-time.
\end{claimdescription}


% =============================
\section{B. Bollig and M. Buttkus, "On the Relative Succinctness of Sentential Decision D...}
% Reference ID: Bollig_2018
% =============================
\begin{claim}
\langfam{SDD}{T} is not quasi-polynomial-time compilable to \langref{FBDD} \citet{Bollig_2018}
\end{claim}
\begin{claimdescription}
The weighted sum function $WS_n$ can be represented by SDDs of size $O(n^3)$ respecting the vtrees used in the construction (Corollary 3 of \citet{Bollig_2018}), but has exponential \langref{FBDD} size. Hence \langfam{SDD}{T} does not compile to \langref{FBDD} in quasi-polynomial size. The union-level lower bound \langref{SDD} $\not\to$ \langref{FBDD} then follows from the polynomial inclusion \langfam{SDD}{T} $\subseteq$ \langref{SDD}.
\end{claimdescription}


% =============================
\section{P. Beame and V. Liew, "New Limits for Knowledge Compilation and Applications to ...}
% Reference ID: Beame_2015
% =============================
\begin{claim}
\langref{FBDD} is not quasi-polynomial-time compilable to \langref{SDD} \citet{Beame_2015}
\end{claim}
\begin{claimdescription}
\citet{Beame_2015} relate the size of SDD representations to best-partition communication complexity and derive simple examples for which FBDDs are exponentially more concise than SDDs. This establishes $P(\langref{FBDD}) \not\subseteq P(\langref{SDD})$.
\end{claimdescription}


% =============================
\section{P. Beame and J. Li and S. Roy and D. Suciu, "Model Counting of Query Expressions...}
% Reference ID: Beame_2013
% =============================
\begin{claim}
\langref{d-DNNF} is not quasi-polynomial-time compilable to \langref{dec-DNNF} \citet{Beame_2013}
\end{claim}
\begin{claimdescription}
Shown by Corollary 3.5 of \citet{Beame_2013}.
\end{claimdescription}


% =============================
\section{K. Pipatsrisawat and A. Darwiche, "New Compilation Languages Based on Structured...}
% Reference ID: Pipatsrisawat_2008
% =============================
\begin{claim}
\langfam{d-SDNNF}{T} is polynomial-time compilable to \langref{d-SDNNF} \citet{Pipatsrisawat_2008}
\end{claim}
\begin{claimdescription}
\langref{d-SDNNF} is the union of the fixed-vtree languages \langfam{d-SDNNF}{T}. Hence any \langfam{d-SDNNF}{T} formula is already a \langref{d-SDNNF} formula, so the compilation is trivial and polynomial-time.
\end{claimdescription}

\begin{claim}
\langfam{SDNNF}{T} is polynomial-time compilable to \langref{SDNNF} \citet{Pipatsrisawat_2008}
\end{claim}
\begin{claimdescription}
\langref{SDNNF} is the union of the fixed-vtree families \langfam{SDNNF}{T}. Hence any \langfam{SDNNF}{T} formula is already an \langref{SDNNF} formula, so the compilation is trivial and polynomial-time.
\end{claimdescription}


% =============================
\section{U. Oztok and A. Darwiche, "On Compiling CNF into Decision-DNNF", International C...}
% Reference ID: Oztok_2014
% =============================
\begin{claim}
\langfam{dec-SDNNF}{T} is polynomial-time compilable to \langref{dec-SDNNF} \citet{Oztok_2014}
\end{claim}
\begin{claimdescription}
\langref{dec-SDNNF} is the union of fixed-vtree \langfam{dec-SDNNF}{T} languages. Hence any fixed-family formula is already a \langref{dec-SDNNF} formula, so the compilation is trivial and polynomial-time.
\end{claimdescription}


% =============================
\section{G. Van den Broeck and A. Darwiche, "On the Role of Canonicity in Knowledge Compi...}
% Reference ID: VanDenBroeck_2015
% =============================
\begin{claim}
\langref{cSDD} is polynomial-time compilable to \langref{SDD} \citet{VanDenBroeck_2015}
\end{claim}
\begin{claimdescription}
As defined by \citet{VanDenBroeck_2015}, \langref{cSDD} is the compressed subset of \langref{SDD}, so compilation is trivial.
\end{claimdescription}

\begin{claim}
\langfam{cSDD}{T} is polynomial-time compilable to \langfam{SDD}{T} \citet{VanDenBroeck_2015}
\end{claim}
\begin{claimdescription}
For any fixed vtree $T$, \langfam{cSDD}{T} is the compressed subset of \langfam{SDD}{T}, so compilation is trivial.
\end{claimdescription}

\begin{claim}
\langref{SDD} is polynomial-time compilable to \langref{d-SDNNF} \citet{VanDenBroeck_2015}
\end{claim}
\begin{claimdescription}
\langref{SDD} is a strict subset of \langref{d-SDNNF}.
\end{claimdescription}

\begin{claim}
\langfam{SDD}{T} is polynomial-time compilable to \langfam{d-SDNNF}{T} \citet{VanDenBroeck_2015}
\end{claim}
\begin{claimdescription}
For any fixed vtree $T$, \langfam{SDD}{T} is a strict subset of \langfam{d-SDNNF}{T}, so compilation is trivial.
\end{claimdescription}


% =============================
\section{A. Darwiche, "SDD: a new canonical representation of propositional knowledge bas...}
% Reference ID: Darwiche_2011
% =============================
\begin{claim}
\langfam{cSDD}{T} is polynomial-time compilable to \langref{cSDD} \citet{Darwiche_2011}
\end{claim}
\begin{claimdescription}
\langref{cSDD} is the union of the fixed-vtree languages \langfam{cSDD}{T}. Hence any \langfam{cSDD}{T} formula is already a \langref{cSDD} formula, so the compilation is trivial and polynomial-time.
\end{claimdescription}

\begin{claim}
\langref{OBDD} is polynomial-time compilable to \langref{cSDD} \citet{Darwiche_2011}
\end{claim}
\begin{claimdescription}
As established by \citet{Darwiche_2011}, every \langref{OBDD} is a right-linear \langref{SDD}, which is already compressed (canonical). Thus, \langref{cSDD} strictly generalizes \langref{OBDD}.
\end{claimdescription}

\begin{claim}
\langfam{OBDD}{<} is polynomial-time compilable to \langfam{cSDD}{T} \citet{Darwiche_2011}
\end{claim}
\begin{claimdescription}
Fix any variable order $<$ and let $T$ be the corresponding right-linear vtree. Every \langfam{OBDD}{<} is a right-linear \langfam{SDD}{T}, and this representation is compressed (canonical), so compilation to \langfam{cSDD}{T} is polynomial-time.
\end{claimdescription}

\begin{claim}
\langfam{SDD}{T} is polynomial-time compilable to \langref{SDD} \citet{Darwiche_2011}
\end{claim}
\begin{claimdescription}
\langref{SDD} is the union of the fixed-vtree languages \langfam{SDD}{T}. Hence any \langfam{SDD}{T} formula is already an \langref{SDD} formula, so the compilation is trivial and polynomial-time.
\end{claimdescription}


% =============================
\section{S. Bova and F. Capelli and S. Mengel and F. Slivovsky, "Knowledge Compilation Me...}
% Reference ID: Bova_2016
% =============================
% separators=Sauerhoff
\begin{claim}
\langref{DNNF} is not quasi-polynomial-time compilable to \langref{d-DNNF} \citet{Bova_2016}
\end{claim}
\begin{claimdescription}
Proposition 7 of \citet{Bova_2016} shows that the Sauerhoff function $S_n$ has \langref{DNNF} size $O(n^2)$. Theorem 9 shows that $S_n$ has \langref{d-DNNF} size $2^{\Omega(n)}$. (Also see \citet{Bova_2014}.)
\end{claimdescription}

% separators=JS
\begin{claim}
\langref{PI} is not quasi-polynomial-time compilable to \langref{DNNF} \citet{Bova_2016}
\end{claim}
\begin{claimdescription}
Proposition 11 of \citet{Bova_2016} shows that the $JS_n$ has \langref{PI} size $O(n^2)$. Proposition 12 of \citet{Bova_2016} shows that $JS_n$ has \langref{DNNF} size $2^{\Omega(n^2)}$. (Also see \citet{Bova_2014}.)
\end{claimdescription}


% =============================
\section{I. Wegener, "The Complexity of Boolean Functions", 1987.}
% Reference ID: Wegener_1987
% =============================
\begin{claim}
\langref{MODS} is polynomial-time compilable to \langref{IP} \citet{Wegener_1987}
\end{claim}
\begin{claimdescription}
The Quine-McCluskey algorithm generates an \langref{IP} formula in polynomial time.
\end{claimdescription}


% =============================
\section{S. Bova, "SDDs Are Exponentially More Succinct than OBDDs", Proceedings of the A...}
% Reference ID: Bova_2016_a
% =============================
% separators=Generalized_HWB
\begin{claim}
\langref{cSDD} is not quasi-polynomial-time compilable to \langref{OBDD} \citet{Bova_2016_a}
\end{claim}
\begin{claimdescription}
Theorem 4 of \citet{Bova_2016_a} constructs a generalized hidden weighted bit function $F_n$ on $2n+1$ variables with \langref{cSDD} size $O(n^3)$ but \langref{OBDD} size $2^{\Omega(n)}$. The function $F_n$ augments $HWB_n$ with fresh variables $y_0, \ldots, y_n$ to ensure the resulting SDD is compressed; conditioning $F_n$ on $y_i = 1$ recovers $HWB_n$, preserving Bryant's exponential \langref{OBDD} lower bound.
\end{claimdescription}


% =============================
\section{H. Vinall-Smeeth, "Structured d-DNNF Is Not Closed Under Negation", Proceedings ...}
% Reference ID: Vinall-Smeeth_2024
% =============================
\begin{claim}
\langref{d-SDNNF} is not polynomial-time compilable to \langref{SDD} \citet{Vinall-Smeeth_2024}
\end{claim}
\begin{claimdescription}
Theorem 1 of \citet{Vinall-Smeeth_2024} proves that there are functions with structured deterministic DNNFs of size $s$ but whose every equivalent \langref{SDD} has size $s^{\widetilde{\Omega}(\log s)}$. The source side is genuinely structured: in the proof of Theorem 2, the hard function is an unambiguous DNF $\psi'$, and the authors fix a vtree $T$ over $\operatorname{var}(\psi')$; each term of $\psi'$ is compiled as a d-DNNF respecting $T$, and their disjunction still respects $T$, with determinism following from unambiguity of $\psi'$. Theorem 1 then transfers this to an SDD lower bound using closure of SDDs under negation.
\end{claimdescription}

\begin{claim}
\langfam{d-SDNNF}{T} is not polynomial-time compilable to \langfam{SDD}{T} \citet{Vinall-Smeeth_2024}
\end{claim}
\begin{claimdescription}
The Vinall-Smeeth construction gives small source circuits that are explicitly structured: for the hard unambiguous DNF $\psi'$, the proof fixes a vtree $T$ over $\operatorname{var}(\psi')$, compiles each term as a d-DNNF respecting $T$, and takes their deterministic disjunction. Thus the source side can be stated as \langfam{d-SDNNF}{T}. The lower bound is against every equivalent \langref{SDD}; since \langfam{SDD}{T} is a subfamily of \langref{SDD}, it also rules out every target vtree family member.
\end{claimdescription}


% =============================
\section{A. Darwiche, "On the tractable counting of theory models and its application to ...}
% Reference ID: Darwiche_2000
% =============================
\begin{claim}
\langref{d-DNNF} is not quasi-polynomial-time compilable to \langref{FBDD} \citet{Darwiche_2000}
\end{claim}
\begin{claimdescription}
\langref{d-DNNF}s can polynomially represent the disjunction of two functions with conflicting variable orders, which forces \langref{FBDD}s to grow exponentially to track the state of both orders within a single "read-once" path.
\end{claimdescription}


% =============================
\section{S. Devadas, "Comparing two-level and ordered binary decision diagram representat...}
% Reference ID: Devadas_1993
% =============================
\begin{claim}
\langref{DNF} is not quasi-polynomial-time compilable to \langref{OBDD} \citet{Devadas_1993}
\end{claim}
\begin{claimdescription}
\citet{Devadas_1993} showed that there exist Boolean functions with polynomial-size \langref{DNF} representations but requiring exponential-size \langref{OBDD} representations for any variable ordering. This establishes $P(\langref{DNF}) \not\subseteq P(\langref{OBDD})$.
\end{claimdescription}


% =============================
\section{F. Capelli, "Structural Restrictions of CNF-formulas: Applications to Model Coun...}
% Reference ID: Capelli_2016
% =============================
\begin{claim}
\langref{FBDD} is not quasi-polynomial-time compilable to \langref{SDNNF} \citet{Capelli_2016,Pipatsrisawat_2010}
\end{claim}
\begin{claimdescription}
Appendix D.2 of \citet{Pipatsrisawat_2010} and Section 6.3 of \citet{Capelli_2016} independently proved this result using different techniques.
\end{claimdescription}


% =============================
\section{J. Gergov and C. Meinel, "Efficient Boolean manipulation with OBDD's can be exte...}
% Reference ID: Gergov_1994
% =============================
% separators=HWB
\begin{claim}
\langref{FBDD} is not quasi-polynomial-time compilable to \langref{OBDD} \citet{Gergov_1994}
\end{claim}
\begin{claimdescription}
The hidden weighted bit function $HWB_n$ has \langref{OBDD} size $2^{\Omega(n/4)}$ \citet{Bryant_1986} but \langref{FBDD} size $O(n^2)$ \citet{Gergov_1994}.
\end{claimdescription}


% =============================
\section{N. S. Kaleyski, "Boolean methods in knowledge compilation", 2016.}
% Reference ID: Kaleyski_2016
% =============================
\begin{claim}
\langref{MODS} is not polynomial-time compilable to \langref{PI} \citet{Kaleyski_2016}
\end{claim}
\begin{claimdescription}
\citet{Kaleyski_2016} constructs a family of Boolean functions $\psi_i$ with \langref{MODS} size $\Theta(2^{2i})$ but \langref{PI} size $\Omega(2^{(i-1)^2/4})$, ruling out a polynomial compilation. Whether an exponential separation exists remains open.
\end{claimdescription}

% =============================
% Bibliography
% =============================
\bibliographystyle{plainnat}
\bibliography{refs}

\end{document}
